% Standalone research entry 204; batch 203--207.
% Original section material: Pasted text(3).txt, source edition 22 September 2026.
% Research additions: 27 September 2026.
% Compile twice with: lualatex 204.tex
% !TeX program = lualatex
% Compile twice: lualatex open_problems_500.tex
% All references and prose are embedded; no BibTeX database or external inputs.
\documentclass[11pt,a4paper]{article}
\usepackage[margin=24mm,headheight=14pt]{geometry}
\usepackage{fontspec}
\setmainfont{Latin Modern Roman}
\setsansfont{Latin Modern Sans}
\setmonofont{Latin Modern Mono}
\usepackage{amsmath,amssymb,mathtools}
\usepackage{unicode-math}
\setmathfont{Latin Modern Math}
\usepackage{microtype}
\usepackage{enumitem,xurl,needspace}
\usepackage[dvipsnames]{xcolor}
\usepackage{hyperref}
\hypersetup{unicode=true,colorlinks=true,linkcolor=MidnightBlue,urlcolor=MidnightBlue,
 pdftitle={500 Mathematical Problem Statements},pdfauthor={Source-annotated compilation},bookmarksnumbered=true}
\usepackage{bookmark}
\usepackage{tocloft}
\setlength{\cftsecnumwidth}{3em}
\cftsetpnumwidth{2.5em}
\cftsetrmarg{4em}
\usepackage{titlesec}
\titleformat{\section}{\Large\bfseries\raggedright}{\thesection}{0.7em}{}
\usepackage{fancyhdr}
\pagestyle{fancy}\fancyhf{}
\fancyhead[L]{\small\sffamily Mathematical problem statements}
\fancyhead[R]{\small Source edition 22}
\fancyfoot[C]{\thepage}
\setlength{\parindent}{0pt}
\setlength{\parskip}{5pt plus 1pt}
\setlength{\emergencystretch}{3em}
\setcounter{tocdepth}{1}
\setcounter{secnumdepth}{1}
\setlist[itemize]{leftmargin=3em,itemsep=5pt,topsep=4pt,parsep=0pt}
\allowdisplaybreaks
\newcommand{\N}{\mathbb{N}}
\newcommand{\Z}{\mathbb{Z}}
\newcommand{\Q}{\mathbb{Q}}
\newcommand{\R}{\mathbb{R}}
\newcommand{\C}{\mathbb{C}}
\newcommand{\F}{\mathbb{F}}
\newcommand{\A}{\mathbb{A}}
\newcommand{\PP}{\mathbb P}
\newcommand{\E}{\mathbb{E}}
\newcommand{\eps}{\varepsilon}
\newcommand{\dd}{\,\mathrm d}
\newcommand{\sref}[2]{\hyperlink{source:#1:#2}{[S#2]}}
\newcommand{\eref}[2]{\hyperlink{extra:#1:#2}{[E#2]}}
% Turn the next switch off for a shorter reading copy. The quotations preserve
% every exactTarget field, including qualifications omitted from a short summary.
\newif\ifshowcatalogue
\showcataloguetrue
\newcommand{\cataloguescope}[1]{\ifshowcatalogue
 \paragraph{Exact scope in the supplied catalogue.}
 \begin{quote}\small #1\end{quote}\fi}

% Standalone wrapper; no external inputs or bibliography files are required.
\fancyhead[L]{\small\sffamily Research notebook: section 204}
\fancyhead[R]{\small 27 September 2026}
\hypersetup{pdftitle={Research attempt 204},pdfauthor={Research notebook}}
\setlength{\parskip}{4pt plus 1pt}
\setlist[itemize]{before=\small\interlinepenalty10000,itemsep=3pt}
\begin{document}
\setcounter{section}{203}
% ================================================================
% problemId: problem.kaplansky-s-zero-divisor-conjecture-for-torsion-free-groups
% source releaseRank: 204; source releaseStatus: open
% exactTarget SHA-256: 77918962cf7a3f1fa1faf31f1f926add306a3b347ddca4c5ea9c04deb86e7c39
\clearpage
\section{\texorpdfstring{Kaplansky's Zero-Divisor Conjecture for Torsion-Free Groups}{Kaplansky's Zero-Divisor Conjecture for Torsion-Free Groups}}\label{204}
\subsection{Definitions and mathematical statement}
For a field $K$ and a group $G$, the group ring $K[G]$ consists of finite sums $\sum_{g\in G}a_gg$, with coefficientwise addition and multiplication determined by the group law and distributivity. A group is torsion-free if $g^m=e$ with $m\ge1$ implies $g=e$. The conjecture is
\[
 a,b\in K[G],\quad ab=0\quad\Longrightarrow\quad a=0\text{ or }b=0
\]
for every torsion-free $G$ and every field $K$. The ring may be noncommutative; ``domain'' here means the displayed absence of nonzero zero-divisors, not commutativity. No finite generation or restriction on the characteristic is imposed.
\par\smallskip\textit{Formulation and concept references:} \sref{204}{1}.
\cataloguescope{Prove that for every torsion-free group G and every field K, the group ring K[G] has no non-trivial zero-divisors, i.e. K[G] is a domain.}
\subsection{Short English statement}
Can two nonzero finite linear combinations of elements of a torsion-free group ever multiply to zero over a field?
% BEGIN RESEARCH ADDITION 204
\subsection{Research attempt}

\paragraph{Current frontier.}
The formulation source develops a geometric search for zero-divisors and
units using product structures satisfying additional conditions that
produce torsion-free CAT$(0)$ groups. Its exclusions and computations
must not be read as support-size theorems for \emph{all} torsion-free
groups.\sref{204}{1}
Gardam's established counterexample concerns the \emph{unit} conjecture:
it supplies a nontrivial unit over $\F_2$, not a pair of nonzero elements
with zero product.\eref{204}{1}
The distinction matters here because nontrivial units are compatible with
a domain. Reduction of the zero-divisor question to finite coefficient
fields is also part of the standard landscape.\eref{204}{2}
No unrestricted zero-divisor counterexample or proof is taken as established
from the sources consulted for this attempt.

\paragraph{Attack route.}
An ordered-group or unique-product argument isolates a surviving product
coefficient, but torsion-freeness alone does not supply that combinatorial
property.\sref{204}{1}
An analytic kernel argument would need extra group-theoretic input and
would not automatically handle positive characteristic. The route chosen
here is instead an explicit counterexample search: separate the finite
coefficient equations from the question whether the required collisions
can occur in a torsion-free group. The latter question is retained as an
exact support-survival condition, not replaced by a check of the
abelianization. A parity obstruction then removes some candidate collision
patterns over every odd characteristic at once.

\paragraph{Attempt.}
\textbf{1. Preserving supports when specializing to a finite field.}
Suppose $\alpha\beta=0$ with both elements nonzero over a field $K$ and a
torsion-free group $G$. Let $R\subset K$ be the subring generated by all
nonzero coefficients of $\alpha,\beta$ and their inverses, together with
the prime subring. It is a finitely generated ring over $\Z$ (or over the
prime field in positive characteristic). Any maximal ideal $\mathfrak m$
of $R$ has finite residue field, as follows.

A field finitely generated as an algebra over $\F_p$ is finite by
Zariski's lemma. A field $F$ finitely generated as a ring over $\Z$ cannot
have characteristic zero: Zariski's lemma would make it a number field,
and its finitely many ring generators would all be integral over
$\Z[1/D]$ for some integer $D>0$. Thus $F$ would be integral over
$\Z[1/D]$. A subring over which a field is integral is itself a field:
the integral equation for $a^{-1}$, multiplied by a suitable power of
$a\ne0$, expresses $a^{-1}$ in that subring. This contradicts the fact
that $\Z[1/D]$ is not a field. This proves the residue-field claim.

Reduction modulo $\mathfrak m$ keeps every listed coefficient nonzero,
because it is a unit of $R$. Distinct elements of $G$ remain distinct.
The specialized factors consequently have \emph{exactly} the original
supports and still multiply to zero. Also, the subgroup generated by
these two finite supports is finitely generated and torsion-free.
This proves, directly, that a counterexample over any field yields one
over some finite field and a finitely generated torsion-free group.
It is a support-preserving version of the familiar reduction, not a new
claim of priority.\eref{204}{2}
It does \emph{not} reduce the problem to $\F_2$, or even to prime fields
without finite extensions.

\textbf{2. An exact collision formulation.}
Normalize the first support elements to $a_1=b_1=e$ by replacing
$(\alpha,\beta)$ by $(a_1^{-1}\alpha,\beta b_1^{-1})$. Write
\[
 \alpha=\sum_{i=1}^m c_i a_i,\qquad
 \beta=\sum_{j=1}^n d_j b_j,
 \quad c_i,d_j\ne0.
\]
Partition the $mn$ cells $(i,j)$ according to equality of $a_i b_j$.
Every block has at least two cells, because a uniquely represented product
would have a nonzero coefficient. No block contains two cells in a
single row or column, by cancellation and distinctness of the supports.
For a proposed partition $\pi$, form the finitely presented group
\[
 U_\pi=\left\langle a_2,\ldots,a_m,b_2,\ldots,b_n\ \middle|\
 a_i b_j=a_kb_l\text{ whenever }(i,j),(k,l)\in C\in\pi
 \right\rangle, \tag{204.1}
\]
with $a_1=b_1=e$. The coefficient equations are the finite system
\[
 \sum_{(i,j)\in C}c_i d_j=0\quad(C\in\pi),
 \qquad \prod_i c_i\prod_jd_j\ne0. \tag{204.2}
\]
Equations (204.2) alone are insufficient: the group (204.1) may have
torsion or identify purportedly distinct support elements.

There is a precise way to keep this obstruction. For any group $U$, set
$U^{(0)}=U$ and let $U^{(r+1)}$ be the quotient of $U^{(r)}$ by the
normal subgroup generated by all its finite-order elements. Define
\[
 T(U)=\varinjlim_{r\ge0}U^{(r)}.
\]
This group is torsion-free. Indeed, if the image of $u$ has finite order
in the direct limit, the relation $u^k=e$ holds at some finite stage;
its image is then killed at the next stage. Moreover every homomorphism
from $U$ into a torsion-free group factors through every $U^{(r)}$,
and hence uniquely through $T(U)$.

\textbf{Lemma 204.1 (support-survival criterion).}
A zero-divisor counterexample exists if and only if there are positive
integers $m,n$, a partition as above, a finite field, and a solution of
(204.2), such that the images of the $a_i$ in $T(U_\pi)$ are pairwise
distinct and the images of the $b_j$ are pairwise distinct.

\emph{Proof.}
For necessity, specialize an actual counterexample as in step 1 and use
its collision partition. The map $U_\pi\to G$ factors through
$T(U_\pi)$, so distinct support elements cannot have collapsed there.
For sufficiency, use the stipulated coefficients and support images in
the group ring of the torsion-free group $T(U_\pi)$. Each block already
has coefficient sum zero. Further product collisions merely combine
zero block sums. The two factors remain nonzero because their individual
supports survive.\hfill$\square$

The construction is an exact logical reduction, not a decision procedure.
The presentation of $U_\pi$ is finite, but the iterated torsion quotient
need not be effectively computable or finitely presented. A large search
with no collapse detected does not establish support survival.

\textbf{3. A coefficient obstruction for pairings.}
Suppose every block has exactly two cells. Orient each pair arbitrarily,
say $((i,j),(k,l))$, and set
\[
 v_C=e_i+f_j-e_k-f_l\in\Z^{m+n}.
\]
Here the $e_i$ and $f_j$ are the row and column basis vectors. A nonzero
coefficient solution requires
\[
 c_i d_j/(c_kd_l)=-1. \tag{204.3}
\]
If integers $t_C$ satisfy $\sum_Ct_Cv_C=0$, multiplication of (204.3)
to these integer powers gives $1=(-1)^{\sum_Ct_C}$.
Thus, in characteristic different from two, every such relation must
have even total coefficient.

\textbf{Lemma 204.2 (exact parity test).}
For a fixed pairing, nonzero coefficients solving its block equations
exist over an algebraic closure of any odd-characteristic field if and
only if
\[
 \sum_C t_Cv_C=0\quad\Longrightarrow\quad
 \sum_Ct_C\equiv0\pmod2
 \quad\text{for all }(t_C)\in\Z^{|\pi|}. \tag{204.4}
\]
When (204.4) holds, a solution exists over a finite extension of each
prescribed odd prime field. In characteristic two, all coefficients
one always solve the pairing equations.

\emph{Proof.}
Necessity was proved above. Let $L$ be the sublattice spanned by the
$v_C$. Under (204.4), the assignment
$\chi(\sum_Ct_Cv_C)=(-1)^{\sum_Ct_C}$ is a well-defined character
$L\to\overline{\F}_p^{\times}$. By Smith normal form, in suitable
integer bases $L$ is generated by $d_1u_1,\ldots,d_ru_r$.
Choose a $d_i$th root of $\chi(d_i u_i)$ for the value of $\chi(u_i)$;
algebraic closure permits this even when $p$ divides $d_i$. Extend by
arbitrary nonzero values on the remaining basis vectors. Evaluating the
extension on the original row and column basis gives the $c_i,d_j$.
Finitely many elements of $\overline{\F}_p$ lie in one finite extension.
The characteristic-two assertion is immediate.\hfill$\square$

This is a coefficient test only. In particular a solution in
$\overline{\F}_p$ cannot repair torsion in $U_\pi$, and a parity
obstruction for pairings says nothing about larger blocks.

\textbf{4. Exhaustive small-pattern experiment and a failed candidate.}
A recursive enumeration paired the lexicographically first unused cell
with each unused cell in a different row and column. No quotient by
relabelling was taken. For each resulting pairing, exact Smith normal
form computed an integer basis for the kernel of the matrix with columns
$v_C$, and checked the parity in (204.4). The finite counts were
\[
\begin{array}{c|r|r}
\text{grid}&\text{labelled admissible pairings}&
\text{pairings excluded in odd characteristic}\\\hline
2\times2&1&0\\
3\times2&2&0\\
3\times4&348&48
\end{array}
\]
These are computational results for this precisely restricted search,
not a proof about arbitrary products in torsion-free groups. The
reproduction script uses integer, not floating-point, normal forms.

One of the $3\times4$ patterns is
\[
\begin{array}{lll}
 C_1:\ (1,1)\sim(2,2),& C_2:\ (1,2)\sim(2,3),&
 C_3:\ (1,3)\sim(3,1),\\
 C_4:\ (1,4)\sim(3,3),& C_5:\ (2,1)\sim(3,4),&
 C_6:\ (2,4)\sim(3,2).
\end{array}
\]
Here $v_{C_2}-v_{C_4}+v_{C_6}=0$, with odd coefficient sum, so no
nonzero coefficient solution exists in odd characteristic. Over $\F_2$
all coefficients one are possible. Does this provide a counterexample?
No: its universal group can be computed exactly. Set $g=a_2$. The
relations successively give
\[
 b_2=g^{-1},\qquad b_3=g^{-2},\qquad a_3=g^{-2},
 \qquad b_4=g^{-4},\qquad g=g^{-6}.
\]
The last remaining relation is automatic. Consequently
$U_\pi\cong\langle g\mid g^7=e\rangle=C_7$. Indeed all the displayed
assignments satisfy the original relations in $C_7$, so no further
collapse was hidden. The formal zero product is
\[
 (1+g+g^5)(1+g^6+g^5+g^3)=0\quad\text{in }\F_2[C_7].
\]
This is a genuine zero product, but its source is torsion. The group
$T(U_\pi)$ is trivial; the supports therefore fail Lemma 204.1.

\textbf{5. A field-independent small-support check.}
A factor supported on two distinct elements cannot be a zero-divisor
in a torsion-free group ring. To see this, normalize it to $c+dg$, where
$c,d\ne0$ and $g\ne e$. If $(c+dg)\beta=0$, then
$\beta=-(d/c)g\beta$. Its finite nonempty support $B$ satisfies
$B=gB$. Thus left multiplication by $g$ permutes $B$. Iterating around
one cycle gives $g^r b=b$, whence $g^r=e$, a contradiction. For a
right factor of support two, the same proof uses right multiplication.
This checks all characteristics and explains why neither a two-term
factor nor the $3\times2$ experiment can produce a torsion-free example.
It is a classical elementary subcase, included here as a check, not as
new progress on the unrestricted conjecture.

\paragraph{Stress test.}
The main possible errors have concrete safeguards. Specialization
inverts \emph{every} nonzero coefficient before reducing. The group
reduction tests distinctness of each support, not distinctness of every
product, since extra product collisions are harmless. Torsion killing
is iterated because a single quotient can create new torsion. The parity
test uses an \emph{integer} kernel basis; an arbitrary rational nullspace
basis without lattice saturation would be inadequate.

The explicit candidate fails in the full universal group, not merely in
its abelianization. Conversely, collapse in an abelianization would not
show collapse in a torsion-free nonabelian group. Neither Gardam's unit
nor the geometric exclusions of the formulation source supply the
missing general support-survival theorem.\eref{204}{1}\sref{204}{1}

\Needspace{10\baselineskip}
\paragraph{Outcome.}
\textbf{Outcome: Conditional reduction.}

The unrestricted question has been reduced to finite-field collision
data together with an exact, potentially non-effective torsion-free
support-survival condition. A complete coefficient criterion for paired
collisions was proved, 348 labelled $3\times4$ patterns were tested, and
one characteristic-two candidate was explicitly traced to $7$-torsion.
No novelty claim is made for the standard finite-field or small-support
reductions; the parity and torsion tests are the deductions developed
in this notebook.

A proof would have to show that every coefficient-feasible pattern
necessarily collapses one of its supports in $T(U_\pi)$, including
unbounded supports and non-pair blocks. A counterexample would require
one rigorously verified surviving pattern. The present experiment
supplies neither, and its finite counts impose no general bound.

% Keep the sources heading with a useful initial bibliography block.
\Needspace{8\baselineskip}
% END RESEARCH ADDITION 204
\subsection{Sources}
\begin{itemize}\raggedright
\item[\textnormal{[S1]}]\hypertarget{source:204:1}{} M. Garg, I. Mineyev, arXiv:2501.07646 (2025).
\url{https://arxiv.org/abs/2501.07646}.
{\small\textit{Catalogue formulation source.}}
% BEGIN REFERENCE ADDITION 204
\item[\textnormal{[E1]}]\hypertarget{extra:204:1}{}
G. Gardam, \emph{A counterexample to the unit conjecture for group rings},
Annals of Mathematics 194 (2021), 967--979,
DOI: 10.4007/annals.2021.194.3.9; main theorem.
\url{https://annals.math.princeton.edu/2021/194-3/p09}.
\item[\textnormal{[E2]}]\hypertarget{extra:204:2}{}
G. Gardam, \emph{Kaplansky's conjectures}, Global Noncommutative Geometry
Seminar, 17 September 2021, slides 3--5, especially the finite-field
reduction on slide 5. The support-preserving ring-specialization proof
above is supplied in the notebook.
\url{https://globalncgseminar.org/wp-content/uploads/2021/09/gncg.pdf}.
% END REFERENCE ADDITION 204
\end{itemize}

\end{document}
